Two households in the same school district, fire district, water utility, and electric
utility service territory are members of the same community in every administratively
meaningful sense.
They pay property taxes together, vote on the same bond measures, and receive services
from the same institutions.
This paper proposes \emph{administrative zone co-membership} as the strongest single
community character signal for redistricting edge weights, and provides a complete
specification of the formula, data pipeline, and legal grounding.

The zone co-membership score is defined as the fraction of available zone types
that two adjacent census tracts share:
$\mathrm{zone\_score}(u,v) = |\{z \in Z : \mathrm{zone}_z(u) = \mathrm{zone}_z(v)\}| \,/\, |Z|$.
This score modifies the METIS edge weight as
$w(u,v) = w_{\mathrm{base}} \times (1 + \alpha \times \mathrm{zone\_score}(u,v))$,
making cuts through administratively co-membered tract pairs geometrically costly.

Six zone types are considered: school district, county subdivision (town/township/borough),
electric utility service territory, fire district, water/sewer district, and police precinct.
School district, county subdivision, and electric utility service territory provide 100\%
national coverage; fire and water districts cover approximately 60\% and 40\% of states,
respectively.
The algorithm degrades gracefully when zone types are unavailable: $|Z|$ is computed
over available types only.

The legal grounding draws on the property tax fiscal bond doctrine established in
\textit{Milliken v.\ Bradley} (1974) and \textit{Reynolds v.\ Sims} (1964),
as well as state utility commission recognition of electric utility rate-payer
communities and the extensive redistricting literature on school district preservation.
Zone co-membership is the most court-defensible community signal in the M-track
because it encodes the administrative relationships that courts have explicitly
recognized as communities of interest, rather than inferring community membership
from economic or geographic similarity.
Empirical validation of the algorithm's effect on NC-14 district plans and New England
town integrity is deferred to post-implementation runs.
